\chapter{CASPT2 class structure and contraction paths}
\label{ch:caspt2-classes}

\section{Purpose and scope}

The perturbation layer of the release is a contracted second-order
treatment of dynamic correlation on a CASSCF reference.  Its energy is
assembled from eight excitation classes, and the class table is part of
the result rather than a diagnostic: the total is meaningful only when
the classes that carry the reference's correlation are present, and the
comparison with an external program is meaningful only when the two
programs use the same class definitions.

The chapter applies to a CASSCF reference whose orbitals are
semicanonical, whose CI vector is state-specific, and whose active space
is small enough for the determinant expansion.  It does not cover a
state-averaged reference, a reference built from a CASCI calculation
alone, or a production-size active space.  The coupled and three-center
paths are value layers; only the diagonal path has a JAX twin.

\section{Theory and derivation}
\label{sec:caspt2-theory}

\subsection{Reference, excitation operators, and the first-order
wavefunction}

Let $|0\rangle$ be the state-specific CASSCF reference,
\begin{equation}
  |0\rangle = \sum_I c_I\,|I\rangle,
  \label{eq:caspt2-reference}
\end{equation}
and let
\begin{equation}
  E_{pq} = \sum_{\sigma\in\{\uparrow,\downarrow\}}
           a^{\dagger}_{p\sigma}a_{q\sigma}
  \label{eq:caspt2-eoperator}
\end{equation}
be the spin-summed excitation operator.  A contracted class function is
a product of two such operators acting on the reference,
\begin{equation}
  |K_p\rangle = O_p|0\rangle,
  \qquad
  O_p = E_{p_1q_1}E_{p_2q_2},
  \label{eq:caspt2-class-op}
\end{equation}
with the index pairs drawn from the core, active, and virtual pools.
The first-order wavefunction is the internally contracted expansion
\begin{equation}
  |1\rangle = \sum_p x_p\,|K_p\rangle,
  \label{eq:caspt2-first-order}
\end{equation}
and the second-order energy is
\begin{equation}
  e_2 = \langle 0|H|1\rangle
      = \sum_p W_p\,x_p,
  \qquad
  W_p = \langle 0|H|K_p\rangle.
  \label{eq:caspt2-e2}
\end{equation}
The functions are orthogonal to the reference for every excitation
class, $\langle 0|K_p\rangle=0$, so no overlap with $|0\rangle$ enters
the energy expression.

\subsection{The projected amplitude equation}

Projecting the Schrödinger equation $(H-E_0)|0+1\rangle=0$ onto the
class space gives
\begin{equation}
  \sum_q \langle K_p|(H_0-E_0)|K_q\rangle\,x_q
  = -\langle K_p|H|0\rangle,
  \label{eq:caspt2-projected}
\end{equation}
or, in matrix form,
\begin{equation}
  (A-E_0 S)\,x = -W,
  \qquad
  A_{pq}=\langle K_p|H_0|K_q\rangle,
  \qquad
  S_{pq}=\langle K_p|K_q\rangle.
  \label{eq:caspt2-system}
\end{equation}
The energy is $e_2=W^{T}x$, and the class decomposition is
$e_{2,\mathrm{class}}=\sum_{p\in\mathrm{class}}W_p x_p$.

Two choices of $H_0$ are implemented.  The diagonal-Fock choice is
\begin{equation}
  H_0^{\mathrm{diag}} = \sum_p \epsilon_p E_{pp},
  \qquad
  E_0 = \sum_p n_p^0\epsilon_p,
  \label{eq:caspt2-h0-diag}
\end{equation}
where $\epsilon_p$ is the semicanonical Fock diagonal and $n_p^0$ is
the reference occupation.  The state-specific generalized-Fock choice
is
\begin{equation}
  H_0 = \sum_{pq} f_{pq}E_{pq},
  \qquad
  f_{pq} = h_{pq}
         + \sum_{rs}D_{rs}\left[(pq|rs)-\tfrac{1}{2}(pr|qs)\right],
  \label{eq:caspt2-h0}
\end{equation}
with $D$ the state 1-RDM and
$E_0=\sum_{pq}f_{pq}D_{pq}$ (the state 1-RDM is symmetric).  The
diagonal of $f$ is the semicanonical
orbital-energy array when the reference orbitals are semicanonical;
its off-diagonal blocks connect determinants of different excitation
classes.  The diagonal choice drops those blocks; the coupled choice
keeps them.

\subsection{The diagonal approximation}

If both $A$ and $S$ are approximated by their diagonal elements, the
projected system decouples:
\begin{equation}
  x_p = -\frac{W_p}{A_{pp}-E_0 S_{pp}},
  \qquad
  e_2 = -\sum_p \frac{W_p^{2}}{A_{pp}-E_0 S_{pp}}.
  \label{eq:caspt2-diagonal}
\end{equation}
For a single contracted function $K$ the quantities are
\begin{equation}
  A_{KK} = \langle K|H_0^{\mathrm{diag}}|K\rangle
         = \sum_p n_p(K)\,\epsilon_p,
  \qquad
  S_{KK} = \langle K|K\rangle,
  \label{eq:caspt2-AK}
\end{equation}
where $n_p(K)$ is the diagonal occupation of orbital $p$ in the
contracted function.  This is the CASPT2-D approximation: the
off-diagonal Fock coupling between functions is neglected.

The diagonal reading is not unique when the class overlap has
degenerate eigenvalues.  The implementation builds the plain-product
basis $E(p_1q_1)E(p_2q_2)|0\rangle$, diagonalizes
\begin{equation}
  S = U\,\mathrm{diag}(\lambda)\,U^{T},
  \label{eq:caspt2-overlap}
\end{equation}
keeps eigenvectors above
$\lambda>\texttt{nthresh}\cdot\lambda_{\max}$, and evaluates
\begin{equation}
  e_2 = -\sum_k \frac{(T^{T}W)_k^{2}}{(T^{T}AT)_{kk}},
  \qquad
  T = U\,\mathrm{diag}(\lambda^{-1/2}),
  \label{eq:caspt2-ortho}
\end{equation}
in the orthonormalized basis.  The option
\texttt{class\_diagonal='block'} instead solves
Eq.~\eqref{eq:caspt2-system} inside the retained subspace,
\begin{equation}
  (A-E_0 S)\,x=-W,
  \qquad
  e_2=W^{T}x,
  \label{eq:caspt2-block}
\end{equation}
which includes the intra-class coupling and is invariant under a change
of basis inside the class.  The option \texttt{'function'} is the
per-function Molcas VP/VM reading of Eq.~\eqref{eq:caspt2-diagonal}.

\subsection{The VP/VM pairing}

For a double excitation with both pairs non-diagonal, the two orderings
of the excitation operators define two components
\begin{equation}
  V_P = \mathrm{cross}+\mathrm{direct},
  \qquad
  V_M = \sqrt{3}\left(\mathrm{cross}-\mathrm{direct}\right),
  \label{eq:caspt2-vpvm}
\end{equation}
which is the OpenMolcas \texttt{mkrhsh} normalization.  The $\sqrt{3}$
multiplies the whole $V_M$ function and cancels in
$W^{2}/(A-E_0S)$ and in the projected linear solve, so the operator
content carried by the code is the relative weight $-1$ between the two
orderings.  Both components enter the energy independently for
off-diagonal pairs and coincide for diagonal pairs.  The $V_M$ component
exists only when both pairs are non-diagonal.

\subsection{The eight classes}

The orbital space is partitioned into core ($C$, \texttt{ncore}
spatial orbitals), active ($A$, \texttt{nact}), and virtual ($V$,
\texttt{nmo}-\texttt{ncore}-\texttt{nact}).  A class is labelled by
the two orbital pools it excites out of and the two it excites into.

\begin{center}
\small
\begin{tabularx}{\textwidth}{>{\raggedright\arraybackslash}p{0.13\textwidth}>{\raggedright\arraybackslash}p{0.14\textwidth}>{\raggedright\arraybackslash}p{0.20\textwidth}X}
\toprule
Class & Pool labels & Operator content & Comment \\
\midrule
\texttt{V0\_ijab} & $(C,C)\to(V,V)$ &
  $E(a,i)E(b,j)$ & two core electrons to two virtuals \\
\texttt{Vm1\_iab} & $(A,C)\to(V,V)$ &
  $E(a,t)E(b,i)$ & one active, one core to two virtuals \\
\texttt{Vm2\_ab} & $(A,A)\to(V,V)$ &
  $E(a,t)E(b,u)$ & two active electrons to two virtuals \\
\texttt{V1\_ija} & $(C,C)\to(A,V)$ &
  $E(t,i)E(a,j)$ & two core electrons to one active and one virtual \\
\texttt{V2\_ij} & $(C,C)\to(A,A)$ &
  $E(t,i)E(u,j)$ & two core electrons to two active orbitals \\
\texttt{V0\_ia} & $(C)\to(V)$, active pair &
  cross $E(v,u)E(a,j)$, direct $E(v,j)E(a,u)$ &
  one inactive hole and one virtual particle dressed by an active-pair
  reorganization; the two components coincide for $u=v$ \\
\texttt{Vm1\_a} & TUVA &
  $E(a,t)E(u,v)$ & three active indices, $a$ virtual,
  $t,u,v$ active \\
\texttt{V1\_i} & ITUV &
  $E(t,i)E(u,v)$ & three active indices, $i$ core,
  $t,u,v$ active \\
\bottomrule
\end{tabularx}
\end{center}

The two three-active-index classes are not single excitations.  They
were single-excitation functions before version 0.3.2; the release
treats them as the plain products shown in the table.

\subsection{The three-active-index classes}
\label{sec:caspt2-tuva}

The products $E(a,t)E(u,v)|0\rangle$ and $E(t,i)E(u,v)|0\rangle$ span a
strongly linearly dependent set.  At \texttt{nthresh}$=10^{-6}$ the
TUVA basis of N$_2$/6-31G CAS(6,6) has rank 1264 of 1728 functions and
the ITUV basis rank 632 of 864, that is, about $73\%$ of the function
count.  A per-function sum $\sum_i W_i^2/(A_{ii}-E_0S_{ii})$ therefore
overcounts the same direction by roughly a factor of three.  Both
classes are evaluated by the orthonormalized diagonal of
Eq.~\eqref{eq:caspt2-ortho} with the deterministic gauge described
above.

On the CAS(2,2) anchors the two classes are below $10^{-9}$~Eh.  On
N$_2$ they contribute $-0.003648$~Eh (\texttt{Vm1\_a}) and
$-0.001625$~Eh (\texttt{V1\_i}), and the default-path total moves from
$0.916$ to $0.989$ of the Wave~12 frozen ORCA CASPT2-D reference.
Against the release reference measured with a single BLAS thread the
default-path ratio is $0.983$; the difference between the two reference
sets is a property of the reference calculation, not of the class
treatment.

A residual remains, and it is a class-pool definition difference rather
than a numerical one.  Read at full core, \texttt{V1\_ija} stands at
$1.553$ of the ORCA CASPT2-D value in every basis.  The diagonal
approximation of that class is gauge dependent in the full-core pool
(random gauges inside the degenerate overlap blocks span $0.879$ to
$0.895$ of the reference, the raw LAPACK basis $1.293$), while its
exact class-block value is gauge invariant at $1.553$.  The ORCA count
of 192 functions is the frozen-core plain-product count,
$m^{2}n_{\mathrm{act}}n_{\mathrm{vir}}=192$ with
$m=n_{\mathrm{core}}-n_{\mathrm{frozen}}=2$; the four-core
plain-product basis of the release reference has 768 functions.  ORCA
freezes the two N $1s$ orbitals, while the release reference and its
class table correlate all four core orbitals, so the two programs count
the same class under different frozen-core settings.  The worked
example below shows that matching the frozen-core setting closes the
class table.

\subsection{The coupled full-Fock system}

The coupled path keeps the off-diagonal blocks of
Eq.~\eqref{eq:caspt2-h0}.  The amplitudes solve
Eq.~\eqref{eq:caspt2-system} with the state-specific $H_0$, and the
class decomposition is the same $e_{2,\mathrm{class}}=\sum W_px_p$.

The overlap matrix is regularly linearly dependent.  It is
diagonalized and eigenvectors with eigenvalue below
$\texttt{rtol}\times\max(1,\lambda_{\max})$ are dropped; the default
\texttt{rtol} is $10^{-8}$.  The assembly costs $n_{\mathrm{fn}}^2$
function-pair contractions plus one $H_0$ application per function.  It
is a NumPy pilot intended for small and medium active spaces, and it
has no JAX twin.  The differentiable production route remains the
diagonal-Fock \texttt{caspt2\_contracted\_e2\_jax} at the CASPT2-D
level.

\subsection{The three-center contraction}

\texttt{core.caspt2\_three\_center} replaces the four-center MO tensor
by a three-center factor
\begin{equation}
  (ab|cd) \approx \sum_P u_{abP}\,u_{cdP},
  \label{eq:caspt2-three-center-factor}
\end{equation}
with the metric factor already folded into $u$, as returned by
\texttt{ri\_eri\_mo} in \texttt{core.ri} with \texttt{full=False}.  To
derive the coupling assembly, expand
\begin{equation}
  g(D) = \sum_J c_J\langle J|H|D\rangle
      = g_h(D) + \sum_{abcd} \eta_{abcd}(D)\,(ab|cd),
  \label{eq:caspt2-g}
\end{equation}
where $\eta(D)$ collects the Slater--Condon coefficients of the ERI
terms.  Substituting Eq.~\eqref{eq:caspt2-three-center-factor},
\begin{equation}
  g(D) = g_h(D)
      + \sum_{abP}\left(\sum_{cd}\eta_{abcd}(D)\,u_{cdP}\right)u_{abP}.
  \label{eq:caspt2-g-factor}
\end{equation}
The inner sum is the pair-density tensor
$Z_{abP}(D)=\sum_{cd}\eta_{abcd}(D)u_{cdP}$, and the coupling of a
contracted function is one contraction
$W=\sum_{abP}Z_{abP}u_{abP}$ plus the one-electron terms.  Neither the
four-center MO tensor nor the spin-orbital tensor is formed: the
working set is $n_{\mathrm{mo}}^2n_{\mathrm{aux}}$ per function,
against $n_{\mathrm{mo}}^4$ for the materialized array and
$16n_{\mathrm{mo}}^4$ for the spin-orbital array of the reference path.

The three-center path is a value layer.  It is not wired into the JAX
chain, and it does not change the approximation level: it computes the
same CASPT2-D energy with the RI factorization of the integrals.

\section{Input and operation}
\label{sec:caspt2-input}

All contraction functions consume the same five objects:

\begin{itemize}
\item \texttt{h1e\_mo}: the one-electron Hamiltonian in the MO basis,
  shape \texttt{(nmo, nmo)}.
\item \texttt{eri\_mo}: the two-electron integrals in chemist notation,
  shape \texttt{(nmo, nmo, nmo, nmo)}; the three-center path replaces
  this object by \texttt{u\_mo}.
\item \texttt{ci}: the state-specific reference CI vector over the
  determinant list.
\item \texttt{dets}: determinant bitmaps with the core spin orbitals
  first, as returned by \texttt{cas\_dets\_for2} in
  \texttt{core.nevpt2.impl} for the arguments \texttt{(ncore,nact,nele)}.
\item \texttt{eps}: the semicanonical Fock diagonal over the same MOs,
  from \texttt{canonicalize\_jax}.  The SCF orbital energies are not a
  substitute: both the CASPT2 and the NEVPT2 denominators are defined
  only in the semicanonical basis.
\end{itemize}

The reference chain is RHF, CASSCF, state-specific CASCI, then
canonicalization with \texttt{canonicalize\_jax}.  The four entry
points are

\begin{center}
\small
\begin{tabularx}{\textwidth}{>{\raggedright\arraybackslash}p{0.44\textwidth}X}
\toprule
Entry point & Role and defaults \\
\midrule
\texttt{caspt2\_contracted\_e2} &
  diagonal-Fock (CASPT2-D level), NumPy; defaults
  \texttt{nthresh}$=10^{-6}$, \texttt{active\_row='replica'},
  \texttt{class\_diagonal='function'}, \texttt{v0\_ia='block'} \\
\texttt{caspt2\_contracted\_e2\_jax} &
  differentiable twin of the diagonal path; same options for
  \texttt{active\_row}, traceable in the CI coefficients \\
\texttt{caspt2\_contracted\_e2\_coupled} &
  state-specific generalized-Fock zeroth order, NumPy; default
  \texttt{rtol}$=10^{-8}$, optional \texttt{fock} and
  \texttt{return\_system} \\
\texttt{caspt2\_contracted\_e2\_three\_center} &
  same eight classes and diagonal denominator, with the coupling from
  a three-center factor; \texttt{u\_mo} shape
  \texttt{(nmo,nmo,naux)} \\
\bottomrule
\end{tabularx}
\end{center}

{\sloppy
The options change how the class space is represented, not what the
class is.  \texttt{active\_row='replica'} drops every $(u,u)$ active row
beyond the first active orbital in \texttt{Vm2\_ab} and \texttt{V2\_ij};
it is exact only for a two-orbital active space and is ignored by the
\texttt{ortho} and \texttt{block} readings, which enumerate every index
tuple.  \texttt{'keep'} retains all rows.  \texttt{class\_diagonal='function'} is the
per-function Molcas pairing, \texttt{'ortho'} the orthonormalized
diagonal of Eq.~\eqref{eq:caspt2-ortho}, and \texttt{'block'} the
projected class solve of Eq.~\eqref{eq:caspt2-block}.
\texttt{v0\_ia='block'} solves one two-component system per
$(u,v,j,a)$; \texttt{'class'} solves the whole class as one projected
system.  The three-active-index classes keep the orthonormalized
diagonal in every setting.\par}

The coupled call is
\begin{lstlisting}[language=Python]
caspt2_contracted_e2_coupled(h1e_mo, eri_mo, ci, dets, *, ncore, nact,
                             nele, fock=None, rtol=1e-8,
                             return_system=False)
\end{lstlisting}
With \texttt{return\_system=True} the function also returns a dictionary
with \texttt{S}, \texttt{A}, \texttt{W}, \texttt{E0}, \texttt{x},
\texttt{T}, \texttt{labels}, \texttt{fock}, \texttt{dm1}, and
\texttt{functions}.  The system is the object to inspect when a class
value is questioned: the smallest eigenvalues of \texttt{S}, the
condition number of the reduced matrix, and the projection of
\texttt{W} on the dropped directions are available without re-running
the assembly.

\section{Worked example}
\label{sec:caspt2-example}

\subsection{N\texorpdfstring{$_2$}{2}/6-31G CAS(6,6): three paths}

The recorded N$_2$/6-31G CAS(6,6) comparison against the ORCA class
table is below.  The three rows are the default per-function reading
(\texttt{class\_diagonal='function'}), the class-block reading
(\texttt{class\_diagonal='block'} with \texttt{v0\_ia='class'}), and
the coupled full-Fock path.
\begin{center}
\small
\begin{tabularx}{\textwidth}{>{\raggedright\arraybackslash}p{0.18\textwidth}>{\raggedright\arraybackslash}p{0.24\textwidth}X}
\toprule
Path & Total ratio & Class evidence \\
\midrule
default (\texttt{function}) &
  $0.989$ (Wave~12 frozen), $0.983$ (release, one BLAS thread) &
  apparent agreement is a cancellation: \texttt{V1\_ija} $1.553$
  ($+55\%$, full core against ORCA's frozen core), \texttt{V2\_ij}
  $0.24$ ($-76\%$), \texttt{V0\_ia} $+10\%$, \texttt{Vm2\_ab} $+12\%$ \\
class-block &
  $1.0232$ &
  \texttt{Vm2\_ab} $1.0000$, \texttt{V0\_ia} $0.9970$; the class
  projections are gauge invariant \\
coupled full-Fock &
  $1.0232$ &
  full-FIC level; diagonal/full coupling ratio $0.9795831$ against the
  ORCA $0.9795829$ \\
\bottomrule
\end{tabularx}
\end{center}
The ratios are the release record in \texttt{README.md} and
\texttt{CHANGELOG.md}; the runnable driver is
\texttt{scripts/caspt2\_n2\_validation.py}, run in three stages:
\begin{lstlisting}
python scripts/caspt2_n2_validation.py --stage ref --outdir out
python scripts/caspt2_n2_validation.py --stage default --outdir out
python scripts/caspt2_n2_validation.py --stage coupled --outdir out
\end{lstlisting}
The default path's total is close to the reference for the wrong
reason: the class residuals have opposite signs and cancel.  The block
reading and the coupled path agree on the total because both include
the intra-class coupling; they differ in what they keep beyond that,
and the agreement is the diagnostic that the class-space projection is
the missing piece of the default path.

{\sloppy
The class table closes once the class pools match ORCA's frozen-core
setting.  With \texttt{frozen\_core}$=2$,
\texttt{class\_diagonal='block'}, \texttt{v0\_ia='class\_ordered'} and
\texttt{three\_active='block'}, all eight classes are within $0.9995$ to
$1.0001$ of the ORCA CASPT2-D values and the total is $-0.071518650$~Eh
against $-0.071520733$~Eh ($0.999971$).  The release defaults keep
\texttt{frozen\_core}$=0$, \texttt{v0\_ia='block'} and
\texttt{three\_active='diagonal'} because the reference and its pins
were generated with those settings; the ORCA-compatible setting is
opt-in.  The default-path total of $0.989$ is therefore a cancellation
between class errors, not an accuracy statement.\par}

\subsection{H\texorpdfstring{$_2$}{2}O/6-31G and CO/6-31G: the
CAS(2,2) class table}

On the two-orbital active space the class-resolved ratios against the
ORCA CASPT2 class table are
\begin{center}
\small
\begin{tabular}{lrrrrr}
\toprule
System & \texttt{V0\_ijab} & \texttt{Vm1\_iab} & \texttt{Vm2\_ab} &
\texttt{V1\_ija} & \texttt{V2\_ij} \\
\midrule
H$_2$O/6-31G & $1.017$ & $1.011$ & $0.999$ & $1.009$ & $0.955$ \\
CO/6-31G    & $1.009$ & $1.011$ & $0.963$ & $0.999$ & $0.987$ \\
\bottomrule
\end{tabular}
\end{center}
The eight-class totals are $1.0103$ (H$_2$O) and $1.0026$ (CO).  The
H$_2$O case uses \texttt{ncore}$=4$ and the CO case \texttt{ncore}$=6$;
both follow the reference chain of the runnable example below.  The
\texttt{V0\_ia} target on H$_2$O is $-0.011135$~Eh (ORCA).  These are
the anchors that establish the class definitions before the
three-active-index treatment is exercised on N$_2$.

\subsection{Running the three contraction routes}

\begin{lstlisting}[language=Python]
import numpy as np
import jax
jax.config.update('jax_enable_x64', True)
import jax.numpy as jnp

from pyscf import gto, scf, mcscf
from diffmrc.core.casci import cs_casci
from diffmrc.core.nevpt2.impl import canonicalize_jax, cas_dets_for2
from diffmrc.core.caspt2 import (caspt2_contracted_e2,
                                 caspt2_contracted_e2_coupled)
from diffmrc.core.caspt2_three_center import (
    caspt2_contracted_e2_three_center)
from diffmrc.core import ri

ncore, nact, nele = 4, 2, 2
nocc = ncore + nact

mol = gto.M(atom='O 0 0 0; H 0.7557 0 0.5880; H -0.7557 0 0.5880',
            basis='6-31g', unit='Angstrom', verbose=0)
mf = scf.RHF(mol).run(verbose=0)
mc = mcscf.CASSCF(mf, nact, nele)
mc.kernel()

C = jnp.asarray(mc.mo_coeff)
h1e = jnp.asarray(mol.intor('int1e_kin') + mol.intor('int1e_nuc'))
eri = jnp.asarray(mol.intor('int2e', aosym='s1'))

def mo_ints(C, h1e, eri):
    h1 = jnp.einsum('ai,ab,bj->ij', C, h1e, C)
    z = jnp.einsum('pi,pqrs->iqrs', C, eri)
    z = jnp.einsum('qj,iqrs->ijrs', C, z)
    z = jnp.einsum('rk,ijrs->ijks', C, z)
    return h1, jnp.einsum('sl,ijks->ijkl', C, z)

nmo = mol.nao_nr()
h1_mo, eri_mo = mo_ints(C, h1e, eri)
_, ci = cs_casci(h1_mo[:nocc, :nocc],
                 eri_mo[:nocc, :nocc, :nocc, :nocc],
                 ncore=ncore, nact=nact, nele_act=nele,
                 nroots=1, return_ci=True)
C2, ci2, eps = canonicalize_jax(C, h1e, eri, ci, ncore=ncore,
                                nact=nact, nele_act=nele)
h1_mo2, eri_mo2 = mo_ints(C2, h1e, eri)
dets = cas_dets_for2(ncore, nact, nele)

e2, breakdown = caspt2_contracted_e2(
    np.asarray(h1_mo2), np.asarray(eri_mo2), np.asarray(ci2), dets,
    ncore=ncore, nact=nact, nele=nele, eps=np.asarray(eps))
for name, value in breakdown.items():
    print(f'{name:10s} {value: .9f}')
print('default', e2)

e2_block, bd_block = caspt2_contracted_e2(
    np.asarray(h1_mo2), np.asarray(eri_mo2), np.asarray(ci2), dets,
    ncore=ncore, nact=nact, nele=nele, eps=np.asarray(eps),
    class_diagonal='block', v0_ia='class')
print('class-block', e2_block)

e2_cpl, bd_cpl, system = caspt2_contracted_e2_coupled(
    np.asarray(h1_mo2), np.asarray(eri_mo2), np.asarray(ci2), dets,
    ncore=ncore, nact=nact, nele=nele, return_system=True)
print('coupled', e2_cpl,
      'min |eig(S)|', np.min(np.abs(np.linalg.eigvalsh(system['S']))))

auxmol, label = ri.make_auxmol(mol)
factors = ri.ri_factors(mol, auxmol)
u_mo, _ = ri.ri_eri_mo(np.asarray(C2), factors.b, factors.jhalf,
                       full=False)
e2_3c, bd_3c = caspt2_contracted_e2_three_center(
    np.asarray(h1_mo2), u_mo, np.asarray(ci2), dets,
    ncore=ncore, nact=nact, nele=nele, eps=np.asarray(eps))
print('three-center', e2_3c,
      'max class difference',
      max(abs(bd_3c[k] - breakdown[k]) for k in breakdown))
\end{lstlisting}

Three points about the call are easy to get wrong.  First,
\texttt{ci} must be the state-specific CI vector of the same determinant
list, not a state-averaged vector and not the natural-orbital CI of a
different active-space ordering.  Second, \texttt{eps} must come from
\texttt{canonicalize\_jax}; passing the RHF orbital energies leaves the
denominators undefined.  Third, \texttt{u\_mo} must be the
metric-folded factor; the raw \texttt{ri.ri\_three\_center} tensor
without \texttt{factors.jhalf} changes the energy by the metric factor.

\section{Pitfalls}
\label{sec:caspt2-limits}

\begin{itemize}
\item \emph{Reading only the total.}  The default N$_2$ total agrees
  with the reference through a cancellation between classes with
  opposite errors.  Read the class table and quote the class ratios
  when the total is used for a comparison.
\item \emph{Comparing a gauge-dependent class.}  When the class overlap
  has degenerate eigenvalues, the diagonal approximation depends on the
  basis chosen inside the block.  Random rotations inside the blocks
  move a class energy by up to $25\%$ (\texttt{V2\_ij}, measured).  Use
  \texttt{class\_diagonal='block'} when the class energy enters a
  decision.
\item \emph{The \texttt{active\_row} replica assumption.}  The default
  \texttt{'replica'} drops $(u,u)$ rows beyond the first active orbital.
  It is exact only for a two-orbital active space; on larger spaces the
  rows are not replicas unless the block projection is used.
\item \emph{The \texttt{V1\_ija} class-space residual.}  Read at full
  core it stands at $1.553$ of the ORCA value in every basis and does
  not shrink with the basis set.  This is a frozen-core definition
  difference, not a convergence artifact: the ORCA count of 192
  functions is the frozen-core plain-product count
  ($m^{2}n_{\mathrm{act}}n_{\mathrm{vir}}=192$ with
  $m=n_{\mathrm{core}}-n_{\mathrm{frozen}}=2$), against $768$ for the
  four-core pool.  Pass \texttt{frozen\_core}$=2$ with
  \texttt{class\_diagonal='block'} to compare this class against ORCA.
\item \emph{Using the three-center path as a different method.}  It is
  the same CASPT2-D energy with RI integrals; the RI fitting error
  enters through \texttt{u\_mo}.
\item \emph{Scaling the coupled path.}  The assembly is
  $n_{\mathrm{fn}}^2$ and is intended for small and medium active
  spaces; it is not a production-scale implementation.
\item \emph{Forgetting the canonical basis.}  The denominators are
  defined in the semicanonical basis.  A CASCI reference with SCF
  orbital energies is not the same calculation.
\end{itemize}

\section{Validation and references}
\label{sec:caspt2-validation}

\begin{center}
\small
\begin{tabularx}{\textwidth}{>{\raggedright\arraybackslash}p{0.22\textwidth}>{\raggedright\arraybackslash}p{0.30\textwidth}X}
\toprule
Check & System & Reference and result \\
\midrule
Diagonal class table &
  H$_2$O/6-31G CAS(2,2) &
  ORCA CASPT2 class table; class ratios $1.017$, $1.011$, $0.999$,
  $1.009$, $0.955$; eight-class total $1.0103$ \\
Diagonal class table &
  CO/6-31G CAS(2,2) &
  class ratios $1.009$, $1.011$, $0.963$, $0.999$, $0.987$; total
  $1.0026$ \\
Three-active-index classes &
  N$_2$/6-31G CAS(6,6) &
  default-path ratio $0.989$ (Wave~12 frozen), $0.983$ (release, one
  BLAS thread); class contributions $-0.003648$~Eh and
  $-0.001625$~Eh \\
Three paths &
  N$_2$/6-31G CAS(6,6) &
  default $0.989$, class-block $1.0232$, coupled $1.0232$;
  \texttt{Vm2\_ab} $1.0000$, \texttt{V0\_ia} $0.9970$,
  \texttt{V1\_ija} $1.553$ (full core), \texttt{V2\_ij} $0.24$ \\
Frozen-core class table &
  N$_2$/6-31G CAS(6,6) &
  \texttt{frozen\_core}$=2$ with the exact class-block reading: eight
  classes $0.9995$-$1.0001$ of ORCA CASPT2-D, total $-0.071518650$~Eh
  against $-0.071520733$~Eh ($0.999971$) \\
Coupled path &
  H$_2$O, CO CAS(2,2) &
  totals $1.0081$, $1.0100$; N$_2$ diagonal/full ratio $0.9795831$
  against $0.9795829$ \\
Three-center parity &
  deterministic exact-factor tensor &
  class-by-class agreement with the materialized path at $10^{-10}$;
  dense-coupling elementwise check \\
\bottomrule
\end{tabularx}
\end{center}

The files behind the table are
\texttt{tests/core/test\_caspt2.py} (class definitions, JAX twin,
option behaviour), \texttt{tests/core/test\_caspt2\_three\_center.py}
(three-center parity and coupling scatter),
\texttt{scripts/caspt2\_n2\_validation.py} (the N$_2$ reference chain
and the ORCA comparison; stages \texttt{ref}, \texttt{default},
\texttt{coupled}), and \texttt{scripts/caspt2\_sparse\_assembly.py}
(sparse class assembly).

The class structure and the internally contracted amplitude equations
follow the OpenMolcas and ORCA conventions documented in
K.~Andersson, P.-A.~Malmqvist, B.~O.~Roos, A.~J.~Sadlej and
K.~Wolinski, \emph{J. Phys. Chem.} \textbf{94}, 5483 (1990);
K.~Andersson, P.-A.~Malmqvist and B.~O.~Roos,
\emph{J. Chem. Phys.} \textbf{96}, 1218 (1992);
C.~Angeli, R.~Cimiraglia, S.~Evangelisti, T.~Leininger and
J.-P.~Malrieu, \emph{J. Chem. Phys.} \textbf{114}, 10252 (2001) and
\textbf{117}, 9138 (2002); and the class table of
D.~Kollmar, K.~Sivalingam and F.~Neese,
\emph{J. Chem. Theory Comput.} (2020), and
D.~Kollmar, Y.~Guo and F.~Neese,
\emph{J. Chem. Theory Comput.} (2021).  The
OpenMolcas implementation records are the \texttt{mkrhsh} and
\texttt{caspt2} sources referred to in the module docstrings.
